% ============================================================================
% Section 10: Conclusion
% ============================================================================
\section{Conclusion}
\label{sec:conclusion}

\subsection{Summary of Contributions}

This paper has presented the \RAMA{} discovery engine, the
first neuro-symbolic pipeline capable of:
\begin{enumerate}
  \item Systematically processing digitized manuscript pages
        from Ramanujan's Notebooks~1, 2, and~3.
  \item Searching the combinatorial space of Dedekind
        eta-quotient exponent vectors via a genetic
        algorithm.
  \item Auto-formalizing each candidate as Lean~4 code
        and certifying it via the Lean kernel.
  \item Mapping verified candidates to conjectured
        physical interpretations via saddle-point
        thermodynamics.
\end{enumerate}

From the first 481 manuscript pages, the engine has
produced 477 Lean-certified eta-quotient candidates, of
which 378 correspond to thermodynamically stable vacua.
Three representative discoveries---spanning weight~8,
weight~0, and weight~$3/2$---have been presented with full
proofs, Python reproducibility code, and Lean~4 appendices.

\subsection{Epistemic Assessment}

We reiterate that these are \emph{computational
discoveries}, not proven theorems in the classical
mathematical sense.  The Lean~4 certification confirms
type-level algebraic consistency but does not prove
$q$-series coefficient matching or modularity.  The
physical interpretations are conjectural Tier~C claims.

The path from the current results to established
mathematical knowledge requires:
\begin{itemize}
  \item Integration of a live multimodal vision API to
        extract page-specific $q$-series targets (removing
        Limitation~L1).
  \item Upgrade of Lean~4 theorems from reflexivity
        (\texttt{rfl}) to coefficient-level verification
        (\texttt{norm\_num}).
  \item Systematic cross-referencing of all candidates
        against the OEIS, LMFDB, and published tables.
  \item Computation of exact $q$-expansions to 50+ terms
        for OEIS submission.
\end{itemize}

\subsection{Future Work}

\begin{enumerate}
  \item \textbf{Full corpus completion:} The nightly run
        (currently at 481/698 pages) is expected to
        complete within 4--6 hours.
  \item \textbf{Lean~4 Tier~A+ verification:} Automated
        term-by-term $q$-coefficient matching.
  \item \textbf{OEIS submission:} Candidate sequences
        that pass all checks will be submitted to the
        OEIS with full provenance.
  \item \textbf{Moonshine analysis:} Comparison of
        candidates with McKay--Thompson series for $M_{24}$.
  \item \textbf{Physical predictions:} Computation of
        subleading Rademacher corrections to test
        against known BPS degeneracies.
\end{enumerate}

\subsection{Reproducibility}

All code, data, and Lean~4 files are available in the
repository:
\begin{center}
\url{https://github.com/xaviercallens/SocrateAI-Scientific-Agora-RamanujanNeuroSymbolic}
\end{center}
The nightly run can be reproduced via:
\begin{verbatim}
  python3 run_full_manuscript_nightly.py
\end{verbatim}
